\section{Result and scope}
This is a continuation of the unresolved Turn~7 existence task. The original
four-dimensional Fabel polynomial in the Zeta workbench admits a literal
arithmetic specialization at the reciprocal cubic of an Erd\H{o}s--Straus
state. Its two channels have different local fibre algebras at the original
prime: the exterior channel supplies ramified quadratic branches, whereas the
middle channel supplies unramified quadratic branches. Both have exactly three
rational local inverse points. The derivative, the full seven-point affine
fibre, its integral specialization and the missing projective sign are retained.

A separate constructive continuation extends the earlier bounded-shape E-to-M
returns to arbitrary negative-four-square shape coordinates on their exact
integer domain. In particular the entire lines $\alpha=-4$ and $\beta=-4$
have original middle returns. Neither result produces the initial state at an
arbitrary prime. No universal ES assertion, completed Turn~7 existence milestone,
historical-priority determination or whole-programme analytic identification is
claimed.

The polynomial and global inverse formulas are antecedents from the Zeta source
\cite[GF1--11, FC32--35]{Zeta}. Its four-dimensional map is the ES reader's
extension, not an additional preimage attributed to the original
three-dimensional Alp\"oge--Fabel example discussed by Tao \cite{Tao}.
The classical Type-I/II coordinates are credited to \cite[\S2]{ET}.

\section{Original arithmetic and the common reciprocal cubic}
Fix a prime $p\equiv1\pmod4$. Retain an original first-half source
\[
 p/4<a<p/2,\quad R=4a-p,\quad u\mid a^2,
 \quad E:R\mid4u+1,\quad M:R\mid u+a.
\]
For $a=\prod_l l^{e_l}$ retain $\beta_l=v_l(u)-e_l\in[-e_l,e_l]$.
With $g=\gcd(a,u)$ the original normalization is
\[
 h=g^2/u,\quad r=u/g,\quad s=a/g,\qquad
 a=hrs,\quad u=hr^2,\quad\gcd(r,s)=1.
\]
For E let $\kappa=(pr+s)/R$ and return
$(x_1,x_2,x_3)=(a,hs\kappa,phr\kappa)$.
For M let $\lambda=(r+s)/R$ and return
$(x_1,x_2,x_3)=(a,phs\lambda,phr\lambda)$.
Multiplication verifies the reciprocal identity. The ordered inverses are
$u=a^2/(Rx_2-pa)$ and $u=pa^2/(Rx_2-pa)$, respectively. The M complement
$u\mapsto a^2/u$ exchanges $r,s$ and the last two denominators; its orientation
is retained rather than suppressed.

Every normalization factor is a $p$-unit. For E use
$pD+1=4hr\kappa$, with $D=(4u+1)/R=(r+\kappa)/s$; in particular $\kappa$
need not be smaller than $p$. For M, $h,r,s<p$ and
$\lambda=(r+s)/R<p$. The three denominators are distinct. In E,
$\kappa>r$ and $pr>s$; in M, $r=s$ would force $r=s=1$ and $R\mid2$.
Their first denominator is below $p/2$ and the other two are above $p/2$.

\begin{lemma}[Original character distinction]\label{lem:character}
For every original divisor $d\mid a^2$, $(d/R)=(d/p)$, with Jacobi symbol on
the left. Every E state has $(h/p)=-1$. Every M state has $(rs/p)=-1$.
\end{lemma}
\begin{proof}
For odd $l\mid a$, quadratic reciprocity and $R\equiv-p\pmod l$ give
$(l/R)=(-1/l)(R/l)=(p/l)=(l/p)$. If $2\mid a$, then $R\equiv-p\pmod8$
and the supplementary law gives $(2/R)=(2/p)$. Multiply over the actual
exponents. At E, $u\equiv-1/4\pmod R$ and $R\equiv3\pmod4$, so
$(u/p)=-1$. Since $u=hr^2$, this proves the first assertion.
At M, $r\equiv-s\pmod R$ and all the factors are units, so
$(r/p)=(r/R)=-(s/R)=-(s/p)$. This proves $(rs/p)=-1$.
These are the classical laws in \cite[\S6.4]{Hatcher}, applied to the retained
original factors.
\end{proof}

Define the labelled reciprocal roots and polynomial
\[
 t_i=p/x_i,\qquad f(T)=\prod_{i=1}^3(T-t_i)
       =T^3-4T^2+\sigma_2T-\sigma_3.                 \tag{1}
\]
Then $2<t_1<4$ and $0<t_2,t_3<2$. The target for the \emph{existing} Fabel
map is
\[
 \mathbf U=(0,-4,\sigma_2,-\sigma_3),\qquad
 H_{\mathbf U}(U,V)=V\prod_i(U-t_iV).               \tag{2}
\]
The leading zero and the coefficient of $U^3V$, namely one, are fixed, not
rescaled away. The inverse on the actual arithmetic domain is $x_i=p/t_i$,
with its positivity and integrality tests and the original E/M divisor inverse.
The coefficient target alone forgets the ordering; $p$ and the labels are kept.

The preceding shape curve has different coordinates:
$\alpha=v-R,\beta=v-D,X=v,Y=2a$, with the additional fixed-prime line
$2Y=p+X-\alpha$. Its inverse is $a=Y/2$, $R=X-\alpha$,
$D=X-\beta$, $u=(RD-1)/4$. Thus the connection to (2) is the explicit
composition through this original state. No identification of two unmarked
cubic curves is asserted.

\section{The literal Fabel map and every inverse branch}
Use $\xi$ for its first source coordinate, distinct from the original $a$.
Fix $i^2=-1$ and put $P=(P_0,P_2,P_3,P_4)$, where
\begin{align*}
 P_0={}&\xi^3z+2\xi^2y-i\xi,\\
 P_2={}&-\xi^3y^2z-2i\xi^2yz+\xi^2w-2\xi^2y^3-10i\xi y^2+3\xi z+y,\\
 P_3={}&2\xi^3y^3z+6i\xi^2y^2z+2\xi^2wy+4\xi^2y^4
             +2i\xi w-4i\xi y^3-2\xi yz+2iz+7y^2,\\
 P_4={}&2\xi^3y^4z+8i\xi^2y^3z+\xi^2wy^2+4\xi^2y^5
             +2i\xi wy+7i\xi y^4-10\xi y^2z-4iyz-w-3y^3.
                                                               \tag{3}
\end{align*}
Its factor coordinates are
\begin{align*}
 b&=i+\xi y,&c&=-i+2\xi y+\xi^2z,\\
 d&=-iy-\xi(iz+2y^2)-\xi^2yz,&e&=2z-7iy^2+\xi w,\\
 f_0&=iw+3iy^3-4yz+\xi(6iy^2z+wy+4y^4)+2\xi^2y^3z.
\end{align*}
Writing $L=\xi U+bV$ and $C=cU^3+dU^2V+eUV^2+f_0V^3$, substitution gives
\[
 \xi d+bc=1,\quad C(-b,\xi)=1,\qquad
 LC=P_0U^4+U^3V+P_2U^2V^2+P_3UV^3+P_4V^4.          \tag{4}
\]
For a simple finite root $t_i$ of (1), let
$d_i=f'(t_i)=(t_i-t_j)(t_i-t_k)$. The full inverse is
\[
 \boxed{\xi^2=d_i^{-1},\quad y=-t_i-i/\xi,\quad
 z=2t_i/\xi+3i/\xi^2,\quad
 w=7it_i^2/\xi+(-4-17t_i)/\xi^2-13i/\xi^3.}         \tag{5}
\]
Both signs of $\xi$ are retained. The finite marked root is recovered by
$t_i=-b/\xi$, and $C(-b,\xi)=\xi^2 f'(t_i)=1$.
Coefficient comparison in (4) determines the quotient factor uniquely; the
factor-coordinate inverse
$y=(b-i)/\xi$, $z=(c+i-2\xi y)/\xi^2$,
$w=(e-2z+7iy^2)/\xi$ then proves (5) and its exhaustiveness.

At $\xi=0$ the single chart inverse is
\[
 q_\infty=(0,-4,\ i(112-\sigma_2)/2,
                 -704+8\sigma_2+\sigma_3).          \tag{6}
\]
It has $b=i,c=-i$. The other signed projective infinity factor has $b=-i$
and is outside this affine chart. Thus (2) has \emph{seven}, not eight,
geometric affine preimages: six finite ones and (6). The missing infinity
sign is different from any arithmetic field or integrality obstruction.

For completeness the constant Jacobian statement used here is verified as
follows. On $\xi\ne0$ use $(\xi,t,A,B)$, with $A=P_0,B=P_2$.
The target equations are
\[
 P_3=\xi^{-2}-4At^3-3t^2-2Bt,\qquad
 P_4=3At^4+2t^3+Bt^2-t\xi^{-2}.
\]
The full source inverse has determinant $-\xi^{-5}$ in those coordinates;
the target determinant is $2\xi^{-5}$. Their ratio is $-2$.
Polynomial continuation proves $\det DP=-2$ everywhere. These are the
antecedent GF10--11 calculations \cite{Zeta}, with their constant and sign
retained; the verifier independently expands the original derivative at every
recorded inverse point.

\section{The complete local algebra distinguishes E from M}
Fix either root $i_p\in\mathbb Q_p$ of $-1$. Its existence follows by lifting
a root modulo $p\equiv1\pmod4$. Normalize $v_p(p)=1$ on all extensions.
The reduced coordinate algebra of the actual fibre is
\[
 \boxed{\mathcal A_{p,x}=\mathbb Q_p\ \times\!
 \prod_{i=1}^3\mathbb Q_p[\xi_i]/(\xi_i^2-d_i^{-1}).}       \tag{7}
\]
The first factor is the single infinity chart. Formula (5) and its inverse
identify the other factors; the derivative is nonzero, so there are no
unrecorded nilpotents. This is the coordinate algebra of the fixed target,
not a newly chosen seven-point model.

\begin{theorem}[Prime-local arithmetic fibre]\label{thm:local}
For every original first-half ES state at a prime $p\equiv1\pmod4$,
$P^{-1}(\mathbf U)$ has exactly three $\mathbb Q_p$-rational points.
More precisely,
\[
 \begin{array}{ll}
 E:&\mathcal A_{p,x}\simeq\mathbb Q_p^3\times K_{\rm ram}\times K_{\rm ram},\\
 M:&\mathcal A_{p,x}\simeq\mathbb Q_p^3\times K_{\rm unr}\times K_{\rm unr},
 \end{array}                                                    \tag{8}
\]
where $K_{\rm ram}/\mathbb Q_p$ is a ramified quadratic extension and
$K_{\rm unr}/\mathbb Q_p$ is the unramified quadratic extension. The repeated
field factors retain their distinct root labels. There is no fully split
seven-point local fibre in the original arithmetic domain.
\end{theorem}
\begin{proof}
For E put $H=hrs\kappa$. Then
\[
 (t_1,t_2,t_3)=H^{-1}(p\kappa,pr,s),
\]
\[
 (d_1,d_2,d_3)=H^{-2}\bigl(
 p(\kappa-r)(p\kappa-s),\
 -p(\kappa-r)(pr-s),\
 (p\kappa-s)(pr-s)\bigr).                              \tag{9}
\]
All factors $h,r,s,\kappa$ are $p$-units. Also $p\nmid\kappa-r$: otherwise
$4hr\kappa\equiv1\pmod p$ would make $4hr^2\equiv1$, contrary to
Lemma~\ref{lem:character}. Hence
\[
 (v_p(d_1),v_p(d_2),v_p(d_3))=(1,1,0),\quad
 d_1/d_2\equiv-1\pmod p,\quad d_3\equiv16\pmod p.       \tag{10}
\]
A unit congruent to a nonzero square modulo the odd prime $p$ is a square in
$\mathbb Q_p$ by Hensel lifting. Thus $d_1/d_2$ is a square and the first
two quadratic factors are copies of one ramified field (their defining
square classes have odd valuation). The last quadratic splits into two
rational points. Together with infinity this proves the E case.

For M put $H=hrs\lambda$. Then
\[
 (t_1,t_2,t_3)=H^{-1}(p\lambda,r,s).
\]
The numerator differences are units: $r,s<p$, $r\ne s$, and $p\nmid r-s$.
Modulo $p$, the derivative classes, up to the common square $H^{-2}$, are
\[
 \boxed{rs,\quad r(r-s),\quad s(s-r).}                  \tag{11}
\]
The first is a nonsquare by Lemma~\ref{lem:character}. Their product is
$-r^2s^2(r-s)^2$, a square since $-1$ is a square modulo $p$. Thus exactly
one of the three classes is a square, and exactly two are nonsquares.
The two nonsquare units have square ratio and define the same unramified
quadratic extension. The remaining quadratic gives two rational points.
This proves (8) and the rational-point count.
\end{proof}

The local facts used in this proof can be made constructive. A root $b_n$
of a unit $z$ modulo $p^n$ has the unique lift
\[
 b_{n+1}=b_n+p^n t,\qquad
 t\equiv-\frac{b_n^2-z}{p^n}(2b_n)^{-1}\pmod p.
\]
Use the compatible rational residue of $z$ to compute the integer quotient.
A nonsquare residue has irreducible separable quadratic reduction, giving the
unramified degree-two extension; a class $p$ times a unit has square root of
valuation $1/2$, giving ramification index two. These are the standard local
field facts in \cite{Milne}; the displayed lift is also checked directly.

An algebra isomorphism over $\mathbb Q_p$ must permute the primitive field
factors. It cannot change their ramification indices. Therefore no such
isomorphism carries the E fibre algebra to the M fibre algebra at the fixed
prime. Equality of the rational-point counts (both three) does not give that
isomorphism. This restricts an \emph{unchanged fibre algebra}; it does not
prohibit an arithmetic map to a different target, constructed below.

\section{Integral specialization and the nonproper boundary}
For E, the exact numerator formula further gives
\[
 v_p(\sigma_2)=1,\quad v_p(\sigma_3)=2,\quad
 v_p(\operatorname{Disc} f)=2,
 \qquad\overline f(T)=T^2(T-4).                         \tag{12}
\]
Indeed $\sigma_2=p\{p\kappa r+s(\kappa+r)\}/H^2$ and
$\kappa+r=sD$ is a unit. If $p\mid D$, then $RD=4u+1$ would make the
nonresidue $u$ equal to the square $-1/4$ modulo $p$, impossible.
The product gives $\sigma_3$ and the three differences give the discriminant.

For each of the four signed branches over $t_1,t_2$, equation (5) has the exact
valuation vector
\[
 \boxed{(v_p(\xi),v_p(y),v_p(z),v_p(w))
                  =(-1/2,1/2,1,1).}                   \tag{13}
\]
The least terms are respectively $\xi$, $-i/\xi$, $3i/\xi^2$ and
$-4/\xi^2$. The coefficient $3$ is a unit since $p\equiv1\pmod4$.
There is no equal-valuation cancellation in these minima.
These are genuine generic-field points but are not integral and have no affine
specialization over the valuation ring. The two branches at $t_3$ and the
infinity branch are integral.

Over an algebraic closure of the residue field, the complete E special fibre
has precisely the following three points:
\[
 (0,-4,56i_p,-704),
 \quad(1/4,-4-4i_p,32+48i_p,-1152-384i_p),
\]
\[
 (-1/4,-4+4i_p,-32+48i_p,-1152+384i_p).                  \tag{14}
\]
The double root zero has no resultant-one factor scaling; the simple root four
has derivative sixteen. Thus the reduction from seven geometric generic points
to three special points is a loss at the affine boundary, not a vanishing source
Jacobian. The determinant $-2$ remains a unit.

For M, the three roots have the distinct reductions $0,r/H,s/H$, and every
derivative value is a unit.  The derivative residues themselves need not be
pairwise distinct: at the valid state
\[
 p=5,\qquad(a,u,R,h,r,s,\lambda,H)=(2,1,3,1,1,2,1,2),
\]
the reduced roots are $(0,3,1)$ whereas the derivative residues are
$(3,1,3)$.  Every geometric inverse in (5) is integral over an unramified
extension. Its geometric special fibre still has seven points; precisely
three are rational over $\mathbb F_p$. Thus counting only rational special
points also fails to distinguish the two channels.

\paragraph{The exact relation to the original conductor.}
The Zeta source constructs fixed invertible complex linear maps $\Phi,\Psi$
with $T_A\Phi=\Psi$ and transported polynomials
$\Phi P\Phi^{-1},\Psi P\Psi^{-1}$ \cite[FC32--35]{Zeta}.
Consequently (5)--(6) transport point by point through that exact complex map,
with its retained original Gram $G$ and norm $q^*Gq$. No Euclidean replacement,
new coefficient fit or linear-conductor kernel is asserted.
The $p$-adic statements here concern the original polynomial over $\mathbb Q(i)$
after the specified embedding in $\mathbb Q_p$. Complex period entries in
$\Psi$ are not thereby declared to have a $\mathbb Q_p$-model. Such a descent
would be an additional theorem.

\paragraph{A nonempty polynomial fibre is not an arithmetic solution.}
The target $(0,-4,0,1)$ has the exact infinity inverse
$(0,-4,56i,-705)$. Its cubic $T^3-4T^2+1$ has no rational root: the rational-root
theorem leaves only $\pm1$, neither a root. Thus it cannot be the reciprocal
cubic of three positive integer denominators at any rational $p$.
The inverse formula alone does not supply the initial arithmetic target.

\section{The exact integral root order and its obligatory denominator}
This is a further arithmetic distinction between the coefficient target and its
split roots. Let $\mathcal O_f=\mathbb Z_p[T]/(f)$, observed by evaluation at the
three original labelled roots. All are integral, and $f$ is monic.

\begin{theorem}[Root-order defect and exact interpolation loss]
For an original E state,
\[
 \boxed{\mathcal O_E\simeq
 \{(z_1,z_2,z_3)\in\mathbb Z_p^3:z_1\equiv z_2\pmod p\}.}       \tag{15a}
\]
Its index in $\mathbb Z_p^3$ is $p$ and its conductor ideal in that ring is
$p\mathbb Z_p\times p\mathbb Z_p\times\mathbb Z_p$.
For M, evaluation gives $\mathcal O_M\simeq\mathbb Z_p^3$.

Given any original E and M state at the same prime, let $x_i$ and $y_i$ be their
ordered reciprocal roots. There is a unique polynomial $F$ of degree at most
two taking $x_i$ to $y_i$, and a unique $G$ of degree at most two taking $y_i$
to $x_i$. Their exact formulas are
\[
 F(T)=\sum_i y_i\frac{(T-x_j)(T-x_k)}{(x_i-x_j)(x_i-x_k)},\qquad
 G(T)=\sum_i x_i\frac{(T-y_j)(T-y_k)}{(y_i-y_j)(y_i-y_k)}.       \tag{15b}
\]
Writing coefficients in increasing degree, $F$ has exact $p$-valuations
$(0,-1,-1)$, while $G\in\mathbb Z_p[T]$.
Moreover $pF\bmod p$ is a nonzero scalar multiple of $T(T-4)$.
The compositions satisfy $G(F(T))=T\pmod{f_E}$ and
$F(G(T))=T\pmod{f_M}$.
\end{theorem}
\begin{proof}
For E, $x_1-x_2$ has valuation one; both differences from $x_3$ are units.
A polynomial over $\mathbb Z_p$ therefore has its first two values congruent
modulo $p$. Conversely, for a tuple with that congruence, interpolation at
$x_1,x_2$ is integral since its divided difference is integral. Its error at
$x_3$ can be corrected by an integral multiple of $(T-x_1)(T-x_2)$, whose
value there is a unit. This proves (15a). The quotient is exactly $\mathbb F_p$
by $z\mapsto z_1-z_2\pmod p$. Multiplication by an arbitrary tuple preserves
the congruence exactly when its first two coordinates are divisible by $p$,
which proves the conductor. All three differences for M are units, giving its
integral interpolation inverse.

Formula (15b) is the original labelled interpolation. The E derivative
valuations imply $pF\in\mathbb Z_p[T]$. It cannot be integral before that
multiplication, since it takes two congruent arguments $x_1,x_2$ to the
incongruent values $y_1,y_2$. In fact modulo $p$, only its second Lagrange
summand survives after multiplication by $p$: $y_1$ is divisible by $p$,
while the third denominator is a unit. Therefore
\[
 pF(T)\equiv (p y_2/d_2)T(T-4)\pmod p,
\]
with a nonzero coefficient. This proves valuation $-1$ for the linear and
quadratic coefficients. The constant term is a unit: its second summand has
valuation zero, and the first and third have positive valuation. The reverse
interpolation has only unit denominators and integral numerators.
Equality of each composition with $T$ at all three distinct labelled roots
proves the two polynomial congruences.
\end{proof}

The exact sequence
\[
 0\longrightarrow\mathcal O_E\longrightarrow\mathbb Z_p^3
 \xrightarrow{\ z\mapsto z_1-z_2\bmod p\ }\mathbb F_p
 \longrightarrow0
\]
has no $\mathbb Z_p$-linear right inverse, because the middle module is
torsion-free. Tensoring with $\mathbb Q_p$ removes this index-$p$ quotient;
that operation loses original integral information. This is the root-order
conductor, explicitly distinct from the Zeta workbench's weighted linear
conductor $T_A$. We do not identify their source modules.

Thus there is a fully explicit rational root-level map between the two cubics,
with a sharp integral obstruction and an integral reverse map. The lack of an
isomorphism between their seven-point local fibre algebras is not described as
unrelatedness. It is the additional quadratic-factor ramification in (8),
which this root interpolation alone does not preserve.

For a small complete illustration, at $p=5$ take the original E denominators
$(2,4,20)$ and M denominators $(2,10,5)$. Their reciprocal root triples are
$(5/2,5/4,1/4)$ and $(5/2,1/2,1)$. The exact maps are
\[
 \boxed{F(T)=\frac{14}{15}T^2-\frac{19}{10}T+\frac{17}{12},\qquad
 G(T)=\frac74T^2-\frac{37}{8}T+\frac{25}{8}.}
\]
The first requires exactly one power of five in its local denominator;
$5F\equiv3T(T-4)\pmod5$. The reverse is integral at five. Both ordered root
triples and their original denominator inverses are retained in the certificate.


\section{An unbounded family of exact same-prime channel returns}
Continue with the original E state, and set
$\alpha=hs^2-R$, $\beta=hs^2-D$. The earlier shape identities give
\[
 h\mid\alpha\beta-1,\qquad p\equiv\alpha\pmod h,
 \qquad p\beta\equiv1\pmod h.                         \tag{15}
\]
\begin{theorem}[Negative-four-square channel transfer]\label{thm:transfer}
Suppose one marked coordinate is $-4c^2$, with $c>0$ integral. Then
$h\equiv3\pmod4$ and $\gcd(h,c)=1$. Put $j=(h+1)/4$.
If $c\mid j$, set
\[
 Q=h,\qquad
 U=\begin{cases}c^2,&\alpha=-4c^2,\\(j/c)^2,&\beta=-4c^2,\end{cases}
 \qquad R_M=(p+4U)/Q,\qquad a_M=(p+R_M)/4.              \tag{16}
\]
Then $U\mid j^2$, $Q\mid p+4U$, $0<R_M<p$, and (16) is an original M state.
With $g_M=\gcd(j,U)$, its complete normalization is
\[
 \boxed{h_M=g_M^2/U,\quad r_M=U/g_M,\quad
 \lambda_M=j/g_M,\quad s_M=R_M\lambda_M-r_M.}           \tag{17}
\]
The new $h_M$ is a square. The ordered return is
$(a_M,ph_Ms_M\lambda_M,ph_Mr_M\lambda_M)$.
Every existing original E state on either shape line $\alpha=-4$ or
$\beta=-4$ is covered, without a bound on its other coordinate.  This does not
assert that either line contains infinitely many source states, or that every
prime supplies a state on one of the lines.
\end{theorem}
\begin{proof}
In the alpha branch, $R=hs^2+4c^2$ and $R\equiv3\pmod4$; in the beta
branch, $D=hs^2+4c^2$ and $RD=4hr^2+1$ with $R\equiv3\pmod4$, so again
$D\equiv3\pmod4$.  In either case $s$ is odd and $h\equiv3\pmod4$.
Equation (15) then gives $h\mid\alpha\beta-1$, which makes $h$ coprime to
$c$.
In the $\alpha$ case the required congruence is immediate.
In the $\beta$ case, $p\equiv-(4c^2)^{-1}\pmod h$ and $4j\equiv1\pmod h$,
so $p+4(j/c)^2\equiv0$. Availability follows from $c\mid j$ in both cases.

For $\alpha$, the original inequality $R=hs^2+4c^2<p$ implies
$4U<p-h<(h-1)p$. For $\beta$, use
$4U\le4j^2=(h+1)^2/4$ and $p>2h$; for $h\ge3$,
$(h+1)^2<8h(h-1)$, so again $4U<(h-1)p$.
Thus $R_M<p$. Its residue is $3\pmod4$, proving the original first-half range.

The valuation normalization of $U\mid j^2$ gives $h_Mr_M\lambda_M=j$,
$h_Mr_M^2=U$, and $\gcd(r_M,\lambda_M)=1$. Also
$a_M=R_Mj-U=h_Mr_Ms_M>0$.
Since $r_M\le j<p$, the equation $QR_M=p+4U$ proves
$\gcd(r_M,R_M)=1$, hence $\gcd(r_M,s_M)=1$.
This is therefore precisely the original M normalization.

For an explicit square-grade formula write \(j=cm\), \(d=\gcd(c,m)\).  Then
\[
 \begin{array}{c|ccc}
   &h_M&r_M&\lambda_M\\ \hline
  \alpha&d^2&c/d&m/d\\
  \beta&d^2&m/d&c/d.
 \end{array}
\]
This also proves that the original shared budget has not been used twice.
\end{proof}

The condition $c\mid j$ is the exact availability domain of the two selected
cofactor words: $c^2\mid j^2$ if and only if $c\mid j$, and $j/c$ must be integral.
It is not a hypothesis inserted into a universal existence theorem. Its failure
can leave other words available, or the entire selected cofactor box empty.
The square-divisor criterion in (16) is classical Type-II arithmetic; what is
proved here is its supply from these actual exterior shapes.

\subsection{Complete inverse fibres and coefficient collisions}
The target retains its cofactor $Q=4h_Mr_M\lambda_M-1$, so it recovers the
original $h$. In the $\alpha$ branch it recovers $c=\sqrt U$; in the $\beta$
branch $c=j/\sqrt U$. The square roots here are positive integers.

For fixed $(p,h,c,\alpha)$, put $N=(p+4c^2)/h$. Every original preimage is
obtained by
\[
 s\mid N,\quad k=N/s,\quad s+k\equiv0\pmod4,\quad r=(s+k)/4,
 \quad a=hrs,\quad u=hr^2,\quad R_s=hs^2+4c^2,
\]
subject to
\[
 \gcd(r,s)=1,\qquad R_s<p,\qquad R_s\mid4hr^2+1.       \tag{17a}
\]
The last divisibility is the original E gate; equivalently
$R_s\mid pr+s$.  Conversely, these tests give
$4hrs-p=R_s$, $u\mid a^2$, and the exact normalized E state.  Indeed, if
$D=(4hr^2+1)/R_s$, then
\[
 pr+s=s(4hr^2+1)-rR_s=R_s(sD-r),
\]
so $\kappa=sD-r$ is integral and
$hs^2-R_s=-4c^2$.  Hence the displayed list, including (17a), is the full
alpha fibre.  The gate cannot be dropped: at
$p=5209,h=95,c=2$, the candidate $s=5,k=11,r=4$ passes the displayed
divisor, parity, gcd and range tests but has
$R_s=2391\nmid6081=4hr^2+1$.

For fixed $(p,h,c,\beta)$, put $F=(4c^2p+1)/h$. For
$1\le s\le\lfloor(p-1)/(2h)\rfloor$, set
\[
 D=hs^2+4c^2,\qquad \Delta_s=s^2(D^2-p)-F.
\]
Keep exactly the positive squares $k^2=\Delta_s$ for which
$sD-k$ is positive and even, and set $r=(sD-k)/2$, $\kappa=sD-r$.
The original coprimality and range tests complete the fibre. This comes from
$4hr(sD-r)=pD+1$ and $k=\kappa-r>0$, so both directions are proved.

Within each marked branch, $c$ is recovered. Equality of the full selected
targets first forces the same retained cofactor
$Q=4h_Mr_M\lambda_M-1=h$.  For this common $h$ and
$j=(h+1)/4$, an $\alpha$ branch with $c_1$ and a $\beta$ branch with $c_2$
select the same $U$ exactly when $c_1c_2=j$. Different original sources in any one fibre remain
separate integer coefficients. Pushing them to one target sums those coefficients;
the kernel is generated by their same-target differences. No injectivity or
one-to-one state count is claimed after those records are discarded.

\subsection{An actual transition, with local ramification changed}
At the hard prime $p=5209$, the original E state is
\[
 (a,R,u,h,r,s,\kappa)=(1330,111,18620,95,14,1,657).
\]
Its shape is $(-16,-576)$, so the two selected branches have $c_1=2,c_2=12$
and $j=24$. Both supply the \emph{same} target $U=4$. Formula (17) gives
\[
 (a_M,R_M,u_M,h_M,r_M,s_M,\lambda_M)=(1316,55,4,4,1,329,6),
\]
\[
 \boxed{\frac4{5209}=\frac1{1316}+\frac1{41130264}+\frac1{125016}.}   \tag{18}
\]
The original E box uses primes $2,5,7,19$. Its M target uses $2,7,47$; the
factor $47$ is part of the reconstructed integer, not an available original
E factor silently adjoined. At this same prime the original Fabel fibre changes
from ramified to unramified by Theorem~\ref{thm:local}, because the target cubic
has changed through an actual arithmetic return. No local algebra isomorphism
between the unchanged fibres is being asserted.

A different hard example is
\[
 p=4561,\quad(a,R,u,h,r,s,\kappa)=(1593,1811,43011,59,27,1,68),
 \quad\beta=-36.
\]
Here $c=3,j=15$ returns
\[
 (a_M,R_M,u_M,h_M,r_M,s_M,\lambda_M)=(1160,79,25,1,5,232,3),
\]
with denominators $(1160,3174456,68415)$.

At the hard prime \(p=3049\), the original E state
\[
 (a,R,u,h,r,s,\kappa)=(806,175,20956,31,26,1,453)
\]
has \(\alpha=-144\), hence \(c=6,j=8\).  The selected \(U=36\) is unavailable
in \(j^2=64\).
More strongly, the entire cofactor box has divisors
\(1,2,4,8,16,32,64\), none satisfying \(31\mid3049+4U\); the required residue
is \(5\pmod{31}\). A different original shell has the middle state
\[
 (a,R,u,h,r,s,\lambda)=(770,31,5,5,1,154,5).
\]
This is an exact local obstruction, not an ES counterexample.

\section{Verification, source distinctions and remaining task}
The release verifier enumerates the complete first-half boxes for all primes
$p\equiv1\pmod4$ through its declared bound, checks every labelled E/M state,
its ordered identity, the local field signatures, and every present negative-square
crossing. It evaluates (3) on all seven exact branches of representative fibres
in the four-dimensional rational algebra generated by $i$ and $\xi$; it does not
assume that this algebra is a field. It computes the original Jacobian by exact
first-order arithmetic and checks the unit-root lifts modulo $p^4$.

The separate program imports none of the main implementation or predecessor
software. It enumerates primitive $h,r,s$ to recover the full state lists,
checks the factor-coordinate product (4), independently enumerates the inverse
crossing fibres, and exhausts the original affine polynomial fibre over
$\mathbb F_{17}$ at both special targets. Execution receipts distinguish these
checks from a mathematical review or a Lean build.

The positive progress is an explicit same-prime return on unbounded shape lines,
and an exact local arithmetic meaning for the inverse branches read in the Zeta
workbench. The root-field signature does not prove that an ES reciprocal target
exists, nor that every remaining exterior shape reaches a middle state.
Turn~7's universal positive-coefficient obligation remains unresolved. We do not
advance to Turn~8 under an assumed occupancy statement.
